% F4 -- What a consecutive-event pair actually is, on two datasets.
\begin{figure}[!htbp]
\centering
\fitwidth{%
\begin{tikzpicture}
\begin{axis}[rep,
  xbar stacked, width=\linewidth, height=3.4cm,
  /pgf/bar width=17pt,
  xmin=0, xmax=100,
  xlabel={share of all consecutive-event pairs (\%)},
  ytick={0,1}, yticklabels={\smob, \sraw},
  y tick label style={align=right},
  grid=none, xmajorgrids,
  nodes near coords, point meta=explicit symbolic,
  every node near coord/.append style={font=\scriptsize,text=white,anchor=center},
  legend style={at={(0.5,-0.70)},anchor=north,legend columns=3,draw=clrule},
  legend cell align=left, legend image code/.code={\draw[#1,draw=none] (0cm,-0.08cm) rectangle (0.28cm,0.14cm);},
]
\addplot[fill=clmodel,draw=none] coordinates {(37,0) [37.0\,\%] (43.6,1) [43.6\,\%]};
\addplot[fill=clink,draw=none]   coordinates {(12,0) [12.0\,\%] (12.4,1) [12.4\,\%]};
\addplot[fill=clmove,draw=none]  coordinates {(51,0) [51.0\,\%] (44.0,1) [44.0\,\%]};
\legend{same antenna name, same mast\,--\,other antenna, the person really moved}
\end{axis}
\end{tikzpicture}}
\caption{The three things a pair of back-to-back records can be}
\label{fig:transition-composition}

\end{figure}
